\section{Scope, conclusion, and claim status}
\label{sec:two-stage-scope}

This note studies the canonical point source $s=e_v$.  The restriction is
substantive: the shifted RPPR load has one positive coordinate, the exact
support is connected to $v$, and the known support-radius theorem is rooted at
one vertex.  General sparse sources remain algebraically meaningful, but the
single-root discovery theorems below do not silently extend to them.

\begin{proposition}[What point-source superposition does and does not give]
\label{prop:two-stage-sparse-source-superposition}
For ordinary linear PPR and a probability source
$s=\sum_{v\in\supp(s)}s_ve_v$,
\begin{equation*}
 x^0(s)=\sum_v s_vx^0(e_v).
\end{equation*}
If a point-source solver costs
$\widetilde O(1/(\sqrt\alpha\,\varepsilon_v))$ and each output has normalized
$\ell_\infty$ error at most $\varepsilon_v$, weighted summation gives error
at most $\sum_vs_v\varepsilon_v$.  The optimal independent allocation yields
\begin{equation}
 \widetilde O\!\left(
  \frac{(\sum_v\sqrt{s_v})^2}
       {\sqrt\alpha\,\epsppr}
 \right)
 \label{eq:two-stage-sparse-source-superposition}
\end{equation}
work, in addition to reading the source.  This factor lies between one and
$\nnz(s)$ and equals $\nnz(s)$ for a uniform sparse source.  It therefore
does not prove the desired additive-$\nnz(s)$ general-source bound.
RPPR itself is nonlinear in $s$, so its supports and discovery traces cannot
be superposed at all.
\end{proposition}

\begin{proof}
Linearity follows from $x^0(s)=Q^{-1}b(s)$.  The triangle inequality gives
the declared combined error.  Minimizing
$\sum_v1/\varepsilon_v$ subject to
$\sum_vs_v\varepsilon_v\leq\epsppr$ gives
$\varepsilon_v\propto1/\sqrt{s_v}$ and
\cref{eq:two-stage-sparse-source-superposition}.  The endpoint bounds follow
from $\sum_vs_v=1$ and Cauchy--Schwarz.  The RPPR obstacle contains the
source-dependent $\ell_1$ active set, which destroys linearity.
\end{proof}

\begin{remark}[A three-vertex exact RPPR obstruction]
On the path $0$--$1$--$2$, take endpoint source weights $1/2$, $\alpha=1/3$,
and $\rho=1/8$.  In degree-unscaled coordinates,
\begin{equation*}
 \widehat H=\frac23D-\frac13A,\qquad
 \widehat\ell=(1/8,-1/12,1/8).
\end{equation*}
Each endpoint component alone has value $3/16$ and contributes only $1/16$
to the middle key, so an independent single-source reporter sees
$1/16-1/12<0$ and stops.  Together the middle key is
$2/16-1/12=1/24>0$ and the middle coordinate must activate.  Thus even the
support-discovery decisions do not compose source by source.
\end{remark}

The main design choice is
\begin{center}
\fbox{screen or solve; discard history; restart on a fixed envelope at most
once.}
\end{center}
This is the reverse of the more delicate accelerated-burn-in narrative.
Acceleration is used only after graph exposure has stopped changing the
operator.  Consequently the main theorem needs no momentum transfer across an
admission, no face-shock convolution, no Perron alignment gate, and no bound
on the number of rejected accelerated windows.

\paragraph{What is proved here.}
\begin{enumerate}[leftmargin=2.2em]
\item A complete point-source set-only lane: one certified Green ball, one
  fixed principal CG solve, and a fair race with APPR.  It is unconditional,
  has a product-scale success branch, and never loses more than a constant
  factor to direct APPR when the cap fails.
\item A certificate-level composition theorem: any Stage~I method returning
  a containing envelope $E$ yields a semantic point-source PPR solver with
  work
  $W_{\rm disc}+\widetilde O(\vol(E)/\sqrt\alpha)$.
\item Exact support is unnecessary.  More strongly, once a containing
  envelope is known, Stage~II can solve the \emph{ordinary linear PPR system}
  on that envelope.  It need not solve RPPR again, project at every step, or
  classify the inactive coordinates inside the envelope.
\item If Stage~I obtained its envelope from a certified numerical RPPR point,
  thresholding that point is already a semantic PPR output.  The fixed linear
  tail is therefore optional and is genuinely needed only by set-only
  discovery engines or by a requested refinement.
\item A fair, interruptible certificate race loses only a constant factor
  relative to the fastest declared discovery engine, while an APPR lane
  guarantees termination.
\item Two AESP-first continuations fit the same portfolio: a signed
  objective-gap handoff to SOR, and a lower-safe mass-completion handoff to
  APPR/SOR.  A new mass-capturing-envelope theorem proves the latter without
  identifying $S^\star$.
\item The companion threshold-batch OP2 theorem, combined with a new
  inner-face energy transfer, gives an arbitrary-graph
  $\widetilde O(1/(\epsppr\sqrt\alpha))$ exact-real randomized algorithm.
  A literal version discards all Stage-I numerical state and solves ordinary
  PPR on the returned inner face; the fastest version returns the certified
  RPPR point directly.
\end{enumerate}

\paragraph{What is imported.}
The RPPR comparison and support-volume facts are source results
\citep{fountoulakis2026complexity}.  The APPR support sandwich, fixed-face
Chebyshev machinery, hard-capped dynamic obstacle protocol, and exact
tree/unicyclic reporters are imported
from the proof-audit provider note \path{aesp_cd_l1_rppr}; their provider
labels are stated at every use.  The AESP-burn-in--SOR handoff is imported
from \path{hybrid_aesp_locsor} and uses the source AESP and SOR frameworks
\citep{huang2025accelerated,chen2023accelerating}.  This note proves the new
modular composition and scheduling statements rather than duplicating those
long proofs.  Section~\ref{sec:two-stage-randomized-op2} imports the
threshold-batch depth and fully charged randomized OP2 theorem from
\path{active_edge_lcp}, then proves the new low-energy-inner-face transfer
needed for a literal ordinary-PPR Stage~II.

\paragraph{What remains open.}
The graph-uniform existence question is closed only in the exact-real
randomized word model of the companion theorem.  Deterministic
finite-precision and coefficient-bit bounds, a practical local randomized
SDD realization, and a persistent face-response implementation remain open.
The mass-envelope rule provides another hard-capped success branch, but the
star barrier below proves that it cannot be uniformly output-volume for an
$\ell_\infty$ target.  APPR remains the unconditional practical baseline and
pays $\Theta(1/(\alpha\rho))$ in the worst case.
